\documentclass[main.tex]{subfiles}
\begin{document}

\begin{quote}
　　古人之观于天地、山川、草木、虫鱼、鸟兽，往往有得，以其求思之深而无不在也。夫夷以近，则游者众；险以远，则至者少。而世之奇伟、瑰怪，非常之观，常在于险远，而人之所罕至焉，故非有志者不能至也。有志矣，不随以止也，然力不足者，亦不能至也。有志与力，而又不随以怠，至于幽暗昏惑而无物以相之，亦不能至也。然力足以至焉，于人为可讥，而在己为有悔；尽吾志也而不能至者，可以无悔矣。
\end{quote}


\begin{flushright}
王安石\\
北宋
\end{flushright}


\begin{quote}
人最宝贵的是生命。生命对于每个人来说只有一次。人的一生应当这样度过：当他回忆往事的时候，不会因为虚度年华而悔恨，也不会因为碌碌无为而羞愧。这样，在临死的时候，他就可以说：“我的整个生命和全部精力，都已经献给了世界上最壮丽的事业--为人类的解放而斗争。” 
\end{quote}


\begin{flushright}
奥斯特洛夫斯基\\
前苏联
\end{flushright}

\end{document}
